\section{Exact finite certificates}
\label{app:certificates}

The following tables provide the finite complement used in
Section~\ref{sec:finite}. Their initial alphabets and maps are defined
in~\eqref{eq:even_finite_alphabet}--\eqref{eq:even_finite_map} and
\eqref{eq:odd_finite_alphabet}--\eqref{eq:odd_finite_map}.
Each cardinality list starts at \(V_0\) and ends at \(V_*\); the
last set is unchanged by one additional application of
\eqref{eq:pruning}. A zero last entry means the stable set is empty.
The coordinate words specify every stable point through
\eqref{eq:word_expansion}, including the doubling convention for the
odd branch. Thus the tables describe point sets, not just cycle counts.

\setcounter{table}{0}
\renewcommand{\thetable}{A.\arabic{table}}
\renewcommand{\theHtable}{appendix.A.\arabic{table}}
\small
\begin{longtable}{@{}r r L{0.43\textwidth}L{0.39\textwidth}@{}}
\caption{Complete even-branch certificates, \(c=-A=0,\ldots,12\).
The words are in normalized integer coordinates.}
\label{tab:even_certificate}\\
\toprule
\(c\) & \(B_c\) & Successive cardinalities & Stable coordinate words \\
\midrule
\endfirsthead
\multicolumn{4}{l}{Table \thetable\ continued.}\\
\toprule
\(c\) & \(B_c\) & Successive cardinalities & Stable coordinate words \\
\midrule
\endhead
\midrule
\multicolumn{4}{r}{Continued on next page.}\\
\endfoot
\bottomrule
\endlastfoot
0 & 2 & \(25,15,10,8,7,6,5,4,3,2\) & \((0)\); \((2)\) \\
1 & 2 & \(25,18,14,12,10,9,8,7\) & \((-1,0,0)\); \((-1,-1,1,1)\) \\
2 & 2 & \(25,17,12,8,7,6\) & \((-2,1,1)\); \((-1,-1,0)\) \\
3 & 3 & \(49,28,18,11,8,6,4,2\) & \((-1)\); \((3)\) \\
4 & 3 & \(49,29,18,14,12,10,8,6\) & \((-2,0)\); \((-2,-2,2,2)\) \\
5 & 3 & \(49,26,16,9,8,7,6\) & \((-3,2,2)\); \((-2,-2,1)\) \\
6 & 3 & \(49,23,12,6,2,0\) & None \\
7 & 3 & \(36,12,6,3,2\) & \((-3,1)\) \\
8 & 4 & \(81,33,14,6,2\) & \((-2)\); \((4)\) \\
9 & 4 & \(64,30,14,12,10,8,6,4\) & \((-3,-3,3,3)\) \\
10 & 4 & \(36,20,12,6\) & \((-4,3,3)\); \((-3,-3,2)\) \\
11 & 4 & \(36,12,6,4,2,0\) & None \\
12 & 4 & \(36,4,2\) & \((-4,2)\) \\
\end{longtable}

\begin{longtable}{@{}r r L{0.43\textwidth}L{0.39\textwidth}@{}}
\caption{Complete odd-branch certificates, \(A=0,-1,\ldots,-16\).
Pruning is carried out in odd doubled coordinates, but the words shown
are in normalized integer coordinates \(x=(z-1)/2\).}
\label{tab:odd_certificate}\\
\toprule
\(A\) & \(B_A\) & Successive cardinalities & Stable coordinate words \\
\midrule
\endfirsthead
\multicolumn{4}{l}{Table \thetable\ continued.}\\
\toprule
\(A\) & \(B_A\) & Successive cardinalities & Stable coordinate words \\
\midrule
\endhead
\midrule
\multicolumn{4}{r}{Continued on next page.}\\
\endfoot
\bottomrule
\endlastfoot
0 & 3 & \(16,8,5,3,2\) & \((0)\); \((1)\) \\
-1 & 4 & \(16,12,9,8,7,6,5,4\) & \((-1,-1,0,0)\) \\
-2 & 5 & \(36,22,17,13,10,8,7,6,5\) & \((-1)\); \((2)\); \((-2,0,0)\) \\
-3 & 5 & \(36,24,17,14,12,10,8,6,5,4\) & \((-2,-2,1,1)\) \\
-4 & 6 & \(36,22,14,8\) & \((-2,-1)\); \((-3,1,1)\); \((-2,-2,0)\) \\
-5 & 6 & \(36,20,12,8,5,2,1,0\) & None \\
-6 & 7 & \(64,32,18,11,8,5,4\) & \((-2)\); \((3)\); \((-3,0)\) \\
-7 & 7 & \(64,32,18,14,10,8,6,4\) & \((-3,-3,2,2)\) \\
-8 & 7 & \(64,28,14,7,6\) & \((-4,2,2)\); \((-3,-3,1)\) \\
-9 & 8 & \(64,24,10,4,2,0\) & None \\
-10 & 8 & \(36,6,2\) & \((-4,1)\) \\
-11 & 8 & \(36,8,2,0\) & None \\
-12 & 9 & \(64,22,10,4,2\) & \((-3)\); \((4)\) \\
-13 & 9 & \(36,20,14,12,10,8,6,4\) & \((-4,-4,3,3)\) \\
-14 & 9 & \(36,20,12,6\) & \((-5,3,3)\); \((-4,-4,2)\) \\
-15 & 9 & \(36,12,6,4,2,0\) & None \\
-16 & 10 & \(36,4,2\) & \((-5,2)\) \\
\end{longtable}
\normalsize

For an independent whole-box reconstruction of these same finite sets,
the even branch can use \([-B_c,B_c]^2\cap\Z^2\) directly. In the odd
branch, the original integer coordinates lie in
\[
 \left[\left\lceil\frac{-B_A-1}{2}\right\rceil,
       \left\lfloor\frac{B_A-1}{2}\right\rfloor\right]^2\cap\Z^2.
\]
For each of these boxes, form the directed adjacency matrix of the
original map \(H_{A,e}\), retaining an edge only when its endpoint
remains in the box. Its Boolean transitive closure has a nonzero
diagonal entry at exactly the periodic vertices. This supplies a
second fully specified finite method for checking the listed complete
sets without the filtered alphabets or a maximum-period guess.
